\section{Conclusion}
\label{sec:conclusion}

This article has described the Research Map as a method for structuring a research plan and examined how its planning relations are represented in an AI-assisted application. The method connects a central problem to objectives, literature and development topics, and a companion objective--method matrix. In the inspected implementation, the agent proposes content within this structure, while the researcher can edit and confirm it. Versioned records, declared links, stage checks, and review states make parts of that process inspectable. Student-authored submissions can additionally pass through an advisor's in-app decision before progression. These correspondences specify what the system supports; they do not establish the scientific adequacy of a plan.

The fictional library example illustrated one possible revision of a research problem and the need to reconsider dependent planning units. It was deliberately constructed rather than run as an application session. Together with the directed literature comparison and source-code inspection, it supports a bounded contribution: a documented account of how this particular planning method is operationalized in a supervised authoring workflow. Related research systems and public applications already offer staged assistance, agentic ideation, revision history, or planning support. The contribution therefore lies in the documented correspondence and its limits, rather than in a claim of exclusive features or demonstrated improvement over those systems.

The present evidence cannot show that the Research Map improves research quality, reduces effort, or works well in routine supervision. Its automated checks are partial, final coherence review is advisory, bibliographic metadata do not verify support for claims, and the constructed case contains no user or advisor observations. The resulting map and its exports should be interpreted as revisable planning artifacts that remain subject to researcher and advisor judgment. Empirical assessment of use and outcomes would require a separate study.
